% GENERATED FILE. Edit canonical lessons, then run npm run generate:books.
% canonical-source-hash: fnv1a64:c8966f4e4ec15bd9
% canonical-lessons: FR-C388-peine, FR-C388-condamner, FR-C388-condamnation, FR-C388-detenu, FR-C388-punir

\chapter{La peine}
\label{ch:fr-la-peine}
\hlchaptermodality{388}

\begin{chapteropening}
\textbf{By the end of this chapter:} \emph{I can talk about punishment: a sentence, condemning someone, a conviction, a prisoner and punishing.}

Complete the last lesson of chapter 388: I can talk about punishment: a sentence, condemning someone, a conviction, a prisoner and punishing.
\end{chapteropening}

\section[la peine]{la peine — the sentence, the penalty; sorrow}

\label{lesson:FR-C388-peine}



\begin{quote}
Before the new one: say the French for guilty, then the French for innocent.
\end{quote}

\subsection*{You'll want to know: la peine}
\textbf{la peine} — \textquotedblleft{}the sentence, the penalty; sorrow\textquotedblright{}.

\begin{quote}
Il a eu une peine de deux ans de prison.
\end{quote}

\begin{quote}
\emph{He got a two-year prison sentence.}
\end{quote}

\emph{Avoir de la peine} — to be sad. \emph{Ce n'est pas la peine} — it is not worth it.

\begin{cousinweb}
\emph{peine} is inherited through Old French from Latin \emph{poena} \textquotedblleft{}a penalty, punishment\textquotedblright{} (Lewis \& Short), itself from Greek \emph{poin\'{\={e}}} \textquotedblleft{}a fine\textquotedblright{}.
\end{cousinweb}

\begin{grammarlens}[title={What you've built}]
One more word for telling, explaining and describing at B1.
\end{grammarlens}

\subsection*{Your turn}
\begin{itemize}
\raggedright
  \item \emph{Say it:} \emph{la peine}
  \item \emph{Say it:} \emph{la peine}, once more
  \item \emph{Say it:} a sentence of your own with \emph{la peine}
  \item \emph{Recall:} say the French for guilty, then the French for innocent, then say \emph{la peine} again
\end{itemize}

\begin{tcolorbox}[breakable,colback=teal!4,colframe=teal!35!black,title={Before you move on}]
What does la peine mean? (\textquotedblleft{}The sentence, the penalty; sorrow\textquotedblright{}.) Say it once more.
\end{tcolorbox}
\section[condamner]{condamner — to sentence, to condemn}

\label{lesson:FR-C388-condamner}



\begin{quote}
Before the new one: say the French for innocent, then the French for the sentence.
\end{quote}

\subsection*{You'll want to know: condamner}
\textbf{condamner} — \textquotedblleft{}to sentence, to condemn\textquotedblright{}.

\begin{quote}
Le tribunal l'a condamné à une amende.
\end{quote}

\begin{quote}
\emph{The court sentenced him to a fine.}
\end{quote}

The \emph{m} is silent: \emph{condaner}.

\begin{cousinweb}
\emph{condamner} is inherited through Old French (EtymDB) from Latin \emph{condemnāre} \textquotedblleft{}to sentence, to condemn\textquotedblright{} (Lewis \& Short), built on \emph{damnāre} \textquotedblleft{}to condemn\textquotedblright{}, from \emph{damnum} \textquotedblleft{}loss, harm\textquotedblright{}.
\end{cousinweb}

\begin{grammarlens}[title={What you've built}]
One more word for telling, explaining and describing at B1.
\end{grammarlens}

\subsection*{Your turn}
\begin{itemize}
\raggedright
  \item \emph{Say it:} \emph{condamner}
  \item \emph{Say it:} \emph{condamner}, once more
  \item \emph{Say it:} a sentence of your own with \emph{condamner}
  \item \emph{Recall:} say the French for innocent, then the French for the sentence, then say \emph{condamner} again
\end{itemize}

\begin{tcolorbox}[breakable,colback=teal!4,colframe=teal!35!black,title={Before you move on}]
What does condamner mean? (\textquotedblleft{}To sentence, to condemn\textquotedblright{}.) Say it once more.
\end{tcolorbox}
\section[la condamnation]{la condamnation — the conviction, the sentence; the condemnation}

\label{lesson:FR-C388-condamnation}



\begin{quote}
Before the new one: say the French for the sentence, then the French for to sentence.
\end{quote}

\subsection*{You'll want to know: la condamnation}
\textbf{la condamnation} — \textquotedblleft{}the conviction, the sentence; the condemnation\textquotedblright{}.

\begin{quote}
Après sa condamnation, il a perdu son emploi.
\end{quote}

\begin{quote}
\emph{After his conviction, he lost his job.}
\end{quote}

The \emph{m} is silent here too.

\begin{cousinweb}
\emph{condamnation} is a learned borrowing from Latin \emph{condemnātiō} \textquotedblleft{}a condemnation\textquotedblright{} (Lewis \& Short), from \emph{condemnāre} \textquotedblleft{}to sentence\textquotedblright{}, built on \emph{damnum} \textquotedblleft{}loss, harm\textquotedblright{}.
\end{cousinweb}

\begin{grammarlens}[title={What you've built}]
One more word for telling, explaining and describing at B1.
\end{grammarlens}

\subsection*{Your turn}
\begin{itemize}
\raggedright
  \item \emph{Say it:} \emph{la condamnation}
  \item \emph{Say it:} \emph{la condamnation}, once more
  \item \emph{Say it:} a sentence of your own with \emph{la condamnation}
  \item \emph{Recall:} say the French for the sentence, then the French for to sentence, then say \emph{la condamnation} again
\end{itemize}

\begin{tcolorbox}[breakable,colback=teal!4,colframe=teal!35!black,title={Before you move on}]
What does la condamnation mean? (\textquotedblleft{}The conviction, the sentence; the condemnation\textquotedblright{}.) Say it once more.
\end{tcolorbox}
\section[le détenu]{le détenu — the prisoner, the inmate}

\label{lesson:FR-C388-detenu}



\begin{quote}
Before the new one: say the French for to sentence, then the French for the conviction.
\end{quote}

\subsection*{You'll want to know: le détenu}
\textbf{le détenu} — \textquotedblleft{}the prisoner, the inmate\textquotedblright{}.

\begin{quote}
Les détenus peuvent recevoir des visites le samedi.
\end{quote}

\begin{quote}
\emph{The inmates can have visitors on Saturdays.}
\end{quote}

Feminine \emph{détenue}. From \emph{détenir}, to hold.

\begin{cousinweb}
\emph{détenu} is the past participle of \emph{détenir} \textquotedblleft{}to hold, to detain\textquotedblright{}, from Latin \emph{dētinēre} \textquotedblleft{}to hold back\textquotedblright{} (Lewis \& Short), built on \emph{tenēre} \textquotedblleft{}to hold\textquotedblright{}.
\end{cousinweb}

\begin{grammarlens}[title={What you've built}]
One more word for telling, explaining and describing at B1.
\end{grammarlens}

\subsection*{Your turn}
\begin{itemize}
\raggedright
  \item \emph{Say it:} \emph{le détenu}
  \item \emph{Say it:} \emph{le détenu}, once more
  \item \emph{Say it:} a sentence of your own with \emph{le détenu}
  \item \emph{Recall:} say the French for to sentence, then the French for the conviction, then say \emph{le détenu} again
\end{itemize}

\begin{tcolorbox}[breakable,colback=teal!4,colframe=teal!35!black,title={Before you move on}]
What does le détenu mean? (\textquotedblleft{}The prisoner, the inmate\textquotedblright{}.) Say it once more.
\end{tcolorbox}
\section[punir]{punir — to punish}

\label{lesson:FR-C388-punir}



\begin{quote}
Before the new one: say the French for the conviction, then the French for the prisoner.
\end{quote}

\subsection*{You'll want to know: punir}
\textbf{punir} — \textquotedblleft{}to punish\textquotedblright{}.

\begin{quote}
Ses parents l'ont puni parce qu'il a cassé la fenêtre.
\end{quote}

\begin{quote}
\emph{His parents punished him because he broke the window.}
\end{quote}

A \emph{-ir} verb like \emph{finir}: \emph{ils punissent}.

\begin{cousinweb}
\emph{punir} comes through Old French \emph{punir} (EtymDB) from Latin \emph{pūnīre} \textquotedblleft{}to inflict a penalty on\textquotedblright{} (Lewis \& Short), from \emph{poena} \textquotedblleft{}penalty\textquotedblright{}.
\end{cousinweb}

\begin{grammarlens}[title={What you've built}]
One more word for telling, explaining and describing at B1.
\end{grammarlens}

\subsection*{Your turn}
\begin{itemize}
\raggedright
  \item \emph{Say it:} \emph{punir}
  \item \emph{Say it:} \emph{punir}, once more
  \item \emph{Say it:} a sentence of your own with \emph{punir}
  \item \emph{Recall:} say the French for the conviction, then the French for the prisoner, then say \emph{punir} again
\end{itemize}

\begin{tcolorbox}[breakable,colback=teal!4,colframe=teal!35!black,title={Before you move on}]
What does punir mean? (\textquotedblleft{}To punish\textquotedblright{}.) Say it once more.
\end{tcolorbox}
